\section{Conclusion}
\label{sec:conclusion}

Bangladesh's July 2024 Internet shutdown provides a high-frequency case study of what happens to digitally mediated production when national connectivity is abruptly removed. Using public GitHub events for 59,158 pre-active accounts across fifteen countries, the paper finds that Bangladesh's daily productive-activity rate fell by 13.26 percentage points relative to a synthetic counterfactual during July 18--23, an 86.8 percent decline. International difference-in-differences, secondary outcomes, placebos, and temporal discontinuity analyses point in the same direction (\cref{tab:main}). A diagnostic event study shows that the aggregate gap narrowed quickly after restoration, although Bangladesh's pre-shutdown weekly estimates were more volatile than those of any donor country.

The commit-level evidence shows that the observed backlog was small relative to normal creation. After the approximately 15:00 UTC July 18 onset, 1,097 commits were authored and subsequently observed publicly by August 10, a 96.35 percent fall against duration-matched pre-shutdown windows (96.75 percent against the original six-day windows). Backlog is about 2.5 percent of normal volume. The original calendar-cohort backlog share of 28.27 percent and onset-cohort share of 68.55 percent describe delayed publication among the surviving observed commits; neither establishes that total software work continued at a comparable rate.

A within-Bangladesh analysis of the restoration period adds a further qualification: reconnection did not immediately restore observable activity. Recovery was faster in Finance- and Information-related repositories than in Education- and Health-related repositories, and activity in Accommodation/Food-related repositories remained depressed throughout the observed window. These sectoral patterns are descriptive comparisons with each sector's own pre-shutdown level rather than causal estimates against an untreated control.

The paper's main limitation is the coincidence of the shutdown with severe domestic unrest. Future work could combine network telemetry at subnational or provider level with developer location at finer geographic resolution, or exploit variation in restoration timing across networks and regions. Such designs could more sharply separate the direct effect of connectivity from the broader social environment. Even with this limitation, the evidence is consistent with observable software-development activity being highly exposed to network availability and with public development platforms providing unusually granular measures of that exposure.
